\begin{abstract}
Airline operations control centers (AOCCs) must prioritize disrupted flight chains with limited review capacity. We ask whether delay rankings reproduce consequence-based shortlists, whether similar shortlists imply similar recovery targets, and how capacity and objective weights relate. Using U.S. airline data from six annual windows (2017--2022), we compare delay- and exposure-based priorities at PRE (before predecessor arrival) and TURN (after arrival, before successor departure), and History+Marginal (HM), History+Point (HP), and History+Joint (HJ) representations under four objectives. At a 10\% review share, six-year equal-weight Delay-only replacement averages are 44.7\% (PRE) and 46.1\% (TURN), with mean normalized screening losses of 0.1397 and 0.1385. Passenger exposure lowers these losses to 0.0985 and 0.0975, while downstream exposure has little effect. The mean screening-loss difference associated with passenger exposure narrows as review share increases, with annual ranges spanning zero at 30\%. HM replaces 1.3--3.8\% of HJ shortlist positions, yet exact recovery-target agreement is 0.581--0.730; HP differs more at both stages. Passenger weighting yields the lowest mean HM--HJ agreement (0.590) and highest HJ activation (0.905); Resource weighting yields the lowest mean activation (0.759). Thus, screening similarity does not imply target agreement; information relevance varies by decision layer, capacity, and objective weights. Results reflect stated references and support populations, not causal effects or realized gains.
\end{abstract}
